\chapter{Lezione 12 --- Discretizzazione dell'impianto, Eulero in avanti e all'indietro, modello esatto con tenuta di ordine zero}

\section{Idea della lezione}

Nelle lezioni precedenti abbiamo seguito soprattutto questa filosofia:
\begin{itemize}
    \item progetto un regolatore continuo \(C(s)\);
    \item poi lo discretizzo ottenendo \(R(z)\), ad esempio con Tustin.
\end{itemize}

In questa lezione compare invece un'altra possibilit\`a:
\begin{itemize}
    \item invece di trasformare il controllore discreto in continuo,
    \item discretizzo direttamente l'impianto continuo,
    \item e ottengo un modello esatto o approssimato dell'impianto nel dominio discreto.
\end{itemize}

Questo approccio \`e particolarmente naturale quando l'impianto \`e descritto in \textbf{forma di stato}.

\section{Segnale in uscita dal convertitore D/A}

Il convertitore digitale-analogico con tenuta di ordine zero produce un segnale
\[
u(t)
\]
\textbf{costante a tratti}.

Quindi:
\begin{itemize}
    \item il segnale \`e continuo a tratti ma non liscio;
    \item presenta discontinuit\`a agli istanti di aggiornamento;
    \item dal punto di vista dell'attuatore, non sempre un ingresso cos\`i \`e desiderabile.
\end{itemize}

In pratica, a seconda dell'attuatore e della struttura fisica del sistema,
un segnale a gradini pu\`o essere accettabile oppure no.

\subsection*{Filtro addizionale di smoothing}

Se il segnale costante a tratti non \`e adatto all'attuatore,
si pu\`o inserire dopo il D/A un ulteriore filtro passa-basso analogico,
in modo da ottenere un comando pi\`u regolare.

L'idea \`e:
\begin{itemize}
    \item il D/A genera un segnale costante a tratti;
    \item il filtro lo \emph{addolcisce};
    \item se il filtro ha banda sufficientemente alta, aggiunge poli ad alta frequenza
    ma non modifica in modo sostanziale le prestazioni in ciclo chiuso nella banda di interesse.
\end{itemize}

Quindi, dal punto di vista progettuale, non serve rifare tutto da capo
se questo filtro \`e collocato abbastanza in alto in frequenza.

\section{Derivata e approssimazioni numeriche}

Negli appunti si richiama il fatto che l'operatore \(s\) nel continuo
pu\`o essere visto come operatore di derivazione.

Se ho una funzione continua \(y(t)\), fissato un passo \(h\),
la derivata si pu\`o approssimare in modo differenziale.

\subsection{Eulero in avanti}

L'approssimazione in avanti \`e
\[
\dot y(t)\approx \frac{y(t+h)-y(t)}{h}.
\]

Nel campionamento con periodo \(T\), questo diventa
\[
\dot y(kT)\approx \frac{y((k+1)T)-y(kT)}{T}.
\]

Applicando la trasformata z:
\[
\mathcal{Z}\left\{\dot y(kT)\right\}
\approx
\frac{zY(z)-Y(z)}{T}
=
\frac{z-1}{T}Y(z).
\]

Da qui si ottiene l'approssimazione
\[
s \approx \frac{z-1}{T},
\]
equivalentemente
\[
z=1+Ts.
\]

\subsection{Osservazione}

Questa \textbf{non} \`e la trasformata di Tustin.
\`E una discretizzazione molto pi\`u rozza, detta \textbf{Eulero in avanti}.

\section{Mappatura del piano \(s\) con Eulero in avanti}

Supponiamo di avere
\[
s=\sigma+j\omega.
\]
Con Eulero in avanti, il corrispondente punto nel piano \(z\) \`e
\[
z=1+Ts = 1+T\sigma + jT\omega.
\]

Questa \`e una mappa affine:
\begin{itemize}
    \item le rette verticali del piano \(s\) restano rette verticali nel piano \(z\);
    \item l'asse immaginario \(\sigma=0\) viene mappato nella retta
    \[
    \Re(z)=1;
    \]
    \item il semipiano sinistro non viene mappato interamente dentro il cerchio unitario.
\end{itemize}

Quindi l'Eulero in avanti \textbf{non preserva automaticamente la stabilit\`a}.

\subsection{Esempio}

Se un polo continuo \`e in
\[
s=-10,
\]
allora il polo discreto approssimato diventa
\[
z=1-10T.
\]

Per avere stabilit\`a discreta occorre
\[
|z|<1
\quad\Longrightarrow\quad
|1-10T|<1.
\]

Da cui segue
\[
0<T<\frac{1}{5}.
\]

Quindi un polo stabile nel continuo pu\`o diventare instabile nel discreto
se \(T\) non \`e sufficientemente piccolo.

\section{Eulero all'indietro}

Si considera ora l'approssimazione all'indietro:
\[
\dot y(kT)\approx \frac{y(kT)-y((k-1)T)}{T}.
\]

Applicando la trasformata z:
\[
\mathcal{Z}\left\{\dot y(kT)\right\}
\approx
\frac{Y(z)-z^{-1}Y(z)}{T}
=
\frac{1-z^{-1}}{T}Y(z).
\]

Quindi
\[
s \approx \frac{1-z^{-1}}{T}
=
\frac{z-1}{zT}.
\]

Risolvendo rispetto a \(z\):
\[
z=\frac{1}{1-sT}.
\]

Questa \`e la discretizzazione di \textbf{Eulero a ritroso} o \textbf{backward Euler}.

\section{Mappatura del piano \(s\) con Eulero all'indietro}

Se
\[
s=\sigma+j\omega,
\]
allora
\[
z=\frac{1}{1-sT}
=
\frac{1}{1-(\sigma+j\omega)T}.
\]

Se \(\sigma<0\), cio\`e il polo continuo sta nel semipiano sinistro,
si verifica che
\[
|z|<1.
\]

Quindi, con Eulero a ritroso:
\begin{itemize}
    \item tutto il semipiano sinistro viene mappato all'interno del cerchio unitario;
    \item i poli stabili nel continuo diventano poli stabili nel discreto.
\end{itemize}

Questo rende l'Eulero a ritroso numericamente molto pi\`u robusto di quello in avanti.

\subsection*{Osservazione pratica}

Negli appunti il messaggio \`e chiaro:
\begin{itemize}
    \item Eulero in avanti pu\`o essere usato come approssimazione veloce e rozza;
    \item Eulero a ritroso \`e pi\`u robusto dal punto di vista della stabilit\`a;
    \item ma se si deve \emph{progettare} un controllore, in generale non conviene usare Eulero:
    \`e meglio usare metodi pi\`u accurati, ad esempio Tustin oppure la discretizzazione esatta dell'impianto.
\end{itemize}

\section{Discretizzare direttamente l'impianto}

La seconda parte della lezione introduce il metodo corretto per ottenere
una rappresentazione discreta dell'impianto a partire da un modello continuo in forma di stato.

Supponiamo che il sistema continuo sia descritto da
\[
\begin{cases}
\dot x(t)=Ax(t)+Bu(t),\\[4pt]
y(t)=Cx(t).
\end{cases}
\]

Sappiamo che la soluzione dell'equazione di stato \`e
\[
x(t)=e^{At}x(0)+\int_0^t e^{A(t-\tau)}Bu(\tau)\,d\tau.
\]

Qui
\[
e^{At}
\]
\`e l'esponenziale di matrice, definito tramite serie:
\[
e^{At}=I+At+\frac{A^2t^2}{2!}+\dots+\frac{(At)^k}{k!}+\dots
\]

Nel caso matriciale, questa serie \`e una \emph{definizione}, non solo una conseguenza
della funzione esponenziale scalare.

\section{Ingresso costante a tratti}

Nel nostro schema di controllo digitale, l'ingresso che entra nell'impianto
\`e costante a tratti a causa della tenuta di ordine zero.

Quindi
\[
u(t)=u_k
\qquad
\forall t\in[kT,(k+1)T).
\]

Questa ipotesi \`e fondamentale:
stiamo discretizzando il sistema continuo tenendo conto del fatto che
l'ingresso viene mantenuto costante tra due campioni successivi.

\section{Passaggio da \(t=kT\) a \(t=(k+1)T\)}

Applichiamo la formula della soluzione tra gli istanti \(kT\) e \((k+1)T\):
\[
x((k+1)T)=e^{AT}x(kT)+\int_{kT}^{(k+1)T} e^{A((k+1)T-\tau)}B\,u(\tau)\,d\tau.
\]

Poich\'e nell'intervallo considerato \(u(\tau)=u_k\), si pu\`o portare fuori dall'integrale:
\[
x((k+1)T)=e^{AT}x(kT)+\left(\int_{kT}^{(k+1)T} e^{A((k+1)T-\tau)}\,d\tau\right)B\,u_k.
\]

Facendo il cambio di variabile
\[
\alpha=(k+1)T-\tau,
\]
si ottiene
\[
x((k+1)T)=e^{AT}x(kT)+\left(\int_0^T e^{A\alpha}\,d\alpha\right)B\,u_k.
\]

Definendo
\[
F=e^{AT},
\qquad
G=\left(\int_0^T e^{A\alpha}\,d\alpha\right)B,
\]
si arriva al modello discreto esatto
\[
x(k+1)=Fx(k)+Gu(k).
\]

Questa \`e l'equazione di stato discreta dell'impianto.

\section{Formula alternativa per \(G\)}

Se \(A\) \`e invertibile, allora si pu\`o anche scrivere
\[
G=A^{-1}(e^{AT}-I)B.
\]

Questa formula \`e utile dal punto di vista computazionale,
ma la forma integrale
\[
G=\int_0^T e^{A\alpha}\,d\alpha\,B
\]
resta la pi\`u generale.

\section{Significato del risultato}

Il modello discreto
\[
x(k+1)=Fx(k)+Gu(k)
\]
non \`e un'approssimazione come l'Eulero in avanti o all'indietro:
\`e il \textbf{modello esatto campionato} del sistema continuo
sotto l'ipotesi di ingresso costante a tratti.

Quindi:
\begin{itemize}
    \item \(F=e^{AT}\) descrive l'evoluzione libera del sistema tra due campioni;
    \item \(G\) descrive l'effetto dell'ingresso costante nell'intervallo di campionamento.
\end{itemize}

Questo \`e il modo corretto di ottenere un impianto discreto equivalente al continuo
quando si ha un D/A con tenuta di ordine zero.

\section{Perch\'e questo metodo \`e importante}

Questa discretizzazione esatta \`e molto importante perch\'e permette:
\begin{itemize}
    \item di trattare l'impianto direttamente in tempo discreto;
    \item di evitare le distorsioni introdotte da approssimazioni rozze della derivata;
    \item di progettare controllori discreti direttamente sul modello campionato dell'impianto.
\end{itemize}

Quindi, invece di fare
\[
C(s)\ \longrightarrow\ R(z),
\]
si pu\`o anche fare
\[
G(s)\ \longrightarrow\ \tilde G(z),
\]
e poi progettare direttamente il regolatore discreto su \(\tilde G(z)\).

\section{Sintesi della lezione}

I punti chiave della lezione sono:
\begin{enumerate}
    \item l'uscita del convertitore D/A \`e un segnale costante a tratti;
    \item se necessario, si pu\`o aggiungere un filtro passa-basso analogico per rendere il comando pi\`u regolare;
    \item l'approssimazione della derivata con Eulero in avanti porta a
    \[
    s\approx \frac{z-1}{T},
    \qquad
    z=1+Ts,
    \]
    ma non preserva sempre la stabilit\`a;
    \item l'approssimazione di Eulero a ritroso porta a
    \[
    s\approx \frac{1-z^{-1}}{T}
    =
    \frac{z-1}{zT},
    \qquad
    z=\frac{1}{1-sT},
    \]
    e mappa il semipiano sinistro dentro il cerchio unitario;
    \item per il progetto del controllo, gli schemi di Eulero non sono in generale la scelta migliore;
    \item se l'impianto \`e in forma di stato
    \[
    \dot x=Ax+Bu,\qquad y=Cx,
    \]
    allora, sotto l'ipotesi di ingresso costante a tratti, il modello discreto esatto \`e
    \[
    x(k+1)=Fx(k)+Gu(k),
    \]
    con
    \[
    F=e^{AT},
    \qquad
    G=\left(\int_0^T e^{A\alpha}\,d\alpha\right)B;
    \]
    \item se \(A\) \`e invertibile, si pu\`o scrivere anche
    \[
    G=A^{-1}(e^{AT}-I)B;
    \]
    \item questa discretizzazione tiene conto in modo esatto della tenuta di ordine zero.
\end{enumerate}